% Simple academic CV (2 pages) for pdflatex
\documentclass[11pt]{article}
\usepackage[margin=1in]{geometry}
\usepackage{titlesec}
\usepackage{enumitem}
\usepackage[hidelinks]{hyperref}
\usepackage{parskip}
\usepackage{setspace}
\usepackage{fancyhdr}

\pagestyle{fancy}
\fancyhf{}
\rhead{Yiwen Ding}
\lhead{Curriculum Vitae}
\cfoot{\thepage}
\setlength{\headheight}{14pt}

\titleformat{\section}{\large\bfseries}{}{0em}{}[\titlerule]
\titleformat{\subsection}{\bfseries}{}{0em}{}
\setlist[itemize]{leftmargin=*, itemsep=0.25em}
\setlength{\parskip}{0.4em}
\setlength{\parindent}{0pt}

\begin{document}
\begin{center}
  {\LARGE \textbf{Yiwen Ding}} \\
  PhD Candidate in Logic, Vrije Universiteit Amsterdam \\
  Email: \href{mailto:dyiwen666@gmail.com}{dyiwen666@gmail.com} \quad
  Phone: +31 (0)6 4734 8246 \\
  GitHub: \href{https://github.com/YiwenDing-Logic}{YiwenDing-Logic} \quad
  Research Portal: \href{https://research.vu.nl/en/persons/yiwen-ding/}{VU Research Profile}
\end{center}

\section*{About}
I am a PhD student in the Department of Ethics, Governance and Society at Vrije Universiteit Amsterdam,
supervised by Alessandra Palmigiano and Alexander Kurz. My research intertwines Category Theory,
Universal Algebra and Coalgebra, and Modal Logic to explore foundational structures that support
reliable reasoning systems, with a focus on non-distributive modal frameworks, causal Kripke semantics,
and neighborhood models.

\section*{Research Interests}
\begin{itemize}
  \item Category Theory
  \item Universal Algebra and Coalgebra
  \item Logic and Artificial Intelligence
  \item Non-distributive Modal Semantics
  \item Causal and Neighborhood Models
\end{itemize}

\section*{Education}
\begin{itemize}
  \item \textbf{Ph.D. Candidate in Logic} \\
  Vrije Universiteit Amsterdam, The Netherlands \hfill Sep 2022 -- Present \\
  Advisor: Prof. Alessandra Palmigiano
  \item \textbf{M.A. in Philosophy} \\
  Shandong University, China \hfill Sep 2019 -- Jun 2022 \\
  Advisor: Prof. Fei Liang
  \item \textbf{B.A. in Law} \\
  Qufu Normal University, China \hfill Sep 2015 -- Jun 2019 \\
  Advisor: Prof. Yingkui Zhang
\end{itemize}

\section*{Selected Publications}
\begin{enumerate}[leftmargin=*, itemsep=0.35em]
  \item Causal Kripke Models. R. Wang, K. Manoorkar, A. Tzimoulis, X. Wang, Yiwen Ding.
  \textit{Synthese} 206(3): 1--36, Article 133 (2025).
  \item Probabilistic Causal Kripke Models. K. Manoorkar, A. Tzimoulis, R. Wang, Yiwen Ding.
  In \textit{Logic, Rationality and Interaction (LORI 2025)}, LNCS 16010.
  \item Monotone Modal Logic Beyond Distributivity. A. Palmigiano, K. Manoorkar, R. Wang,
  Yiwen Ding, A. Kurz. In \textit{Indian Conference on Logic and Its Applications 2025}, LNCS 15402.
  \item Toward the van Benthem Characterization Theorem for Non-Distributive Modal Logic.
  K. Manoorkar, M. Panettiere, R. Wang, Yiwen Ding. \textit{Advances in Modal Logic 2024}, Prague.
  \item Causal Kripke Models. K. Manoorkar, A. Tzimoulis, R. Wang, X. Wang, Yiwen Ding.
  \textit{Electronic Proceedings in Theoretical Computer Science} 379: 185--200 (2025).
  \item Fuzzy Lattice-based Description Logic. K. Manoorkar, Yiwen Ding.
  \textit{Electronic Proceedings in Theoretical Computer Science} 421: 44--63 (2025).
\end{enumerate}

\section*{Conference Presentations}
\begin{itemize}
  \item Logic, Rationality and Interaction 2025, Xi'an, China
  \item Indian Conference on Logic and Its Applications 2025
  \item Advances in Modal Logic 2024, Prague, Czech Republic
  \item Topology, Algebra, and Categories in Logic 2024, Barcelona, Spain
  \item 7th Asian Workshop on Philosophical Logic 2025, Kolkata, India
\end{itemize}

\section*{Academic Visits}
\begin{itemize}
  \item University of the Witwatersrand, South Africa -- Jan 2023
  \item Chapman University, USA -- Apr 2023
  \item Chapman University, USA -- May 2024
\end{itemize}

\section*{Scholarships \& Support}
\begin{itemize}
  \item China Scholarship Council full funding for PhD studies in the Netherlands (2022--present).
\end{itemize}

\vfill
\begin{center}
  \small Updated: March 11, 2026
\end{center}

\end{document}
